% ════════════════════════════════════════════════════════════════════
% SECTION 1 — INTRODUCTION
% ════════════════════════════════════════════════════════════════════
\section{Introduction}\label{sec:intro}

Solar-flare prediction from photospheric magnetogram time series is a
benchmark-driven machine-learning problem. The community's public
reference is SWAN-SF~\cite{angryk2020}, and the machine-learning
literature on it has grown around a small number of evaluation
conventions: randomised splits of pooled windows, the cleaned and
rebalanced export, and a reporting culture built on ROC/PR summary
statistics. Those conventions carry risk. The benchmark's own creators
warned that sliding-window temporal coherence ``invalidates the
randomness assumption, thus impacting all sampling practices''~\cite{angryk2019};
subsequent work has echoed the warning~\cite{ahmadzadeh2021}, and the
cleaned export's release paper explicitly prescribes pairing
temporally-preceding training partitions with a later test
partition~\cite{eskandari2024}. Machine-learning flare-prediction
models built on photospheric vector-magnetogram
summaries---the SHARP parameter suite~\cite{bobra2015} and its
descendants---have demonstrated statistically skilful classification of
flaring versus non-flaring active regions~\cite{leka2019,georgoulis2021};
the predictive signal lives in quantities such as total unsigned
current helicity, magnetic free energy, and flux-emergence proxies. What
has not existed, per the documented search protocol of
Section~\ref{sec:related}, is a federated treatment of SWAN-SF---with
its cleaned and raw substrates, its leakage structure, and its
centralised baselines---under an evaluation protocol that audits its own
benchmark artefacts.

Why federate a public benchmark whose data is, in reality, freely
poolable? Two honest reasons. Methodologically, cross-silo federated
learning~\cite{mcmahan2017,kairouz2021} is the standard instrument for
studying non-IID data heterogeneity under a controlled protocol: it
prescribes exactly what moves between data holders (model parameters)
and exactly what does not (raw observations), which makes it a clean
lens for asking how much of a benchmark's apparent difficulty survives
when the data is deliberately fragmented. Operationally, the
SHARP-style parameters at the centre of this pipeline are dominated by
a single public instrument (SDO/HMI via JSOC); the six-regional-agency
narrative sometimes attached to federated space-weather proposals does
not describe this data and is not claimed here. The six clients of
this study are a \emph{simulation protocol}---Dirichlet label skew
across disjoint shards---not a claim about any institution, and the
present study does not claim that any agency's flare data cannot be
pooled.

In cross-silo FL, each data-holder trains a shared architecture locally
and transmits model parameters to an aggregation server, which averages
them into a new global model; raw measurements do not move. FL has
proven transformative in medicine and digital health, where privacy
constraints really are binding~\cite{rieke2020}. Federated flare
forecasting itself is not new: Fu et al.\ posted a federated
transfer-learning preprint in July 2023 (Section~\ref{sec:related})
\cite{fu2023}. As a body of practice, FL also carries its own
evaluation-protocol hazards---statistics that travel with weights
(most visibly BatchNorm running moments, a failure mode documented
since FedBN~\cite{li2021fedbn}), drift monitors that saturate, and
checkpoint rules that quietly restore the wrong model---and a benchmark
audit is exactly the place to expose how those hazards interact with a
published comparison.

This paper provides that treatment, and its character is an audit: we
ask what a federated pipeline actually achieves on SWAN-SF, what the
benchmark's own artefacts do to the answer, and which parts of the
observed behaviour belong to the algorithms rather than to the
evaluation machinery. We present a complete, reproducible federated
framework---implemented, versioned, and released as open
source---that \emph{simulates} six clients (labelled A--F; no real
institution participates), each holding a non-IID shard of the
benchmark in both its cleaned/rebalanced form and its raw, naturally
imbalanced form, and evaluates whether federation can approach the
centralised upper bound without any data movement. The framework emits
M/X flare probabilities whose calibration, thresholds, and
false-alarm rates are themselves audited quantities in this paper---the
evaluation finds they do not survive the prevalence shift, and no
operational alerting claim is made. The downstream
flare$\to$CME$\to$storm$\to$GIC chain is physically far from the
quantity modelled here and is not evaluated
(Section~\ref{sec:limitations}; see \cite{boteler2019} for that
literature). The specific contributions are:

\begin{enumerate}[leftmargin=2em, itemsep=3pt]
  \item \textbf{A problem formulation and system architecture} that maps
  flare-probability estimation onto a cross-silo federated
  classification problem with a safety-critical, recall-weighted
  objective, and that operationalises client heterogeneity as a
  concrete, simulable protocol (Section~\ref{sec:problem} and
  Section~\ref{sec:method}).
  \item \textbf{A non-IID federation protocol for SWAN-SF} that combines
  Dirichlet label-skew partitioning across six simulated regional
  clients, a server-aware Fed-Focal loss \cite{sarkar2020} for extreme
  class imbalance, and optional distribution-aware aggregation and
  per-client SMOTE oversampling (both studied as ablation arms; SMOTE
  is inert on the cleaned variant and an executed negative at natural
  prevalence), on top of FedAvg
  \cite{mcmahan2017}, FedProx \cite{li2020}, and SCAFFOLD
  \cite{karimireddy2020} training loops
  (Section~\ref{sec:method}).
  \item \textbf{A frozen, leakage-audited evaluation protocol and a
  data-provenance audit of the benchmark export.} The protocol
  comprises an immutable train/validation/test contract, calibration
  and threshold selection on validation only, single-pass test
  evaluation, a centralised-MLP architecture control, five-seed
  statistics, a 24-cell component-ablation grid, region-disjoint
  split generation, natural-prevalence operating points, untouched
  client holdouts, a training-budget audit, event-level alerting
  metrics, and a calibration-comparison harness
  (Section~\ref{sec:experiments}). The provenance audit aligns every
  cleaned window to the raw SWAN-SF instances (Harvard Dataverse;
  normalization-invariant matching) and finds that the cleaned
  export's same-partition train/test pairing shares instances---including 100\% of the flaring test windows and 85.7--90\%
  TimeGAN-synthetic training positives---so the study re-evaluates all
  headline claims on a leakage-free temporally-preceding fold
  (Section~\ref{sec:provenance}). The leakage phenomenon itself is
  documented prior art---the benchmark's creators warned about
  random-split overlap in 2019~\cite{angryk2019} and the cleaned
  export's release prescribes temporally-preceding
  folds~\cite{eskandari2024}; what is new here is the instance-level
  quantification, performed inside a federated study.
  \item \textbf{A leakage-free federated evaluation.} On the
  temporally-preceding fold (training partitions 1--4, single-pass
  test on partition 5, natural 1.31\% prevalence), all arms
  generalise (ROC-AUC 0.906--0.978); the in-partition
  federated-versus-centralised gap disappears (FedProx 0.976/0.308
  vs architecture-matched centralised MLP 0.973/0.382); the
  threshold-transfer and calibration failures persist; and an
  event-level evaluation on 65 physical M/X-class events shows
  detection is not the deployability boundary---XGBoost's false-alarm
  burden is roughly $7\times$ lower than every neural arm's
  (Section~\ref{sec:res-leakfree}).
  \item \textbf{A raw-substrate re-execution and the encoder-conditional
  collapse.} Retraining the identical frozen protocol on the raw,
  unbalanced benchmark (natural prevalence, no synthetic positives, the
  release's imputation and normalisation reproduced in-pipeline with
  train-only parameters) separates the arms decisively: every
  federated \emph{MLP} arm degrades sharply (FedAvg detects zero of 65
  events; SCAFFOLD 0.768/0.057) while the same protocol with the
  sequential SolarLSTM restores the federated standing (FedProx-LSTM
  0.970/0.404; SCAFFOLD-LSTM the substrate's highest event-level
  detection, 55/65 at 0.4 false-alarm windows per day --- an
  implementation-case result, since the arm runs as the sign-inverted,
  cold-start variant disclosed in Section~\ref{sec:method}), a second seed
  holds all four sequence arms within $\pm$0.005 ROC-AUC, and
  per-client SMOTE at natural prevalence is a clean negative---so the
  balanced-substrate standing is a property of the flattened-feature
  MLP, not of federation per se (Sections~\ref{sec:res-raw}--\ref{sec:res-lstmrep}).
  \item \textbf{A BatchNorm-transport diagnostic and a controlled
  no-BatchNorm experiment, which re-attribute the in-partition
  FedAvg-versus-FedProx gap.} The underlying mechanism---weight
  averaging without normalisation statistics---is established
  knowledge in the FL literature since FedBN~\cite{li2021fedbn} and
  its successors~\cite{wang2023bn,guerraoui2024bn,bnscaffold2024};
  the contribution here is the composite failure-mode dissection on
  this pipeline and the self-correction of our own earlier artefact.
  The shipped implementation transports
  model parameters only, so BatchNorm running statistics remain at
  initialisation on the server for the entire run; under that
  evaluation the FedAvg global model drifts from 0.950 to 0.098 test
  ROC-AUC over 50 rounds, and a saturated validation monitor ($F_1
  = 2p/(1{+}p)$ whenever every prediction is positive) restores the
  collapsed round-20 checkpoint as the ``best'' model. The same
  weights score 0.951 under correct normalisation and 0.953 even at
  round 50; with BatchNorm layers removed, both FedAvg and FedProx
  train stably to near-identical test quality, so the published
  algorithm gap is substantially an evaluation-implementation
  artefact (Sections~\ref{sec:res-bn}--\ref{sec:res-nobn}).
  \item \textbf{An in-partition data-locality study, disclosed as
  such.} On the same-partition protocol (331{,}185 test windows,
  1.88\% prevalence) whose instance overlap the audit quantified,
  FedProx attains ROC-AUC 0.954 (97.7\% of the XGBoost reference's
  0.977) against the architecture-matched centralised MLP's 0.888, a
  comparison whose optimisation budgets are documented and
  budget-matched; these numbers are retained as protocol-stability
  evidence with the corrected attribution above, not as
  generalisation estimates (Section~\ref{sec:results}).
  \item \textbf{A cross-model interpretability audit.} SHAP analysis
  of both the centralised XGBoost reference and the federated global
  model, with physical-parameter-group aggregation, shows both
  concentrate decisions on current helicity, the Schrijver
  flux-emergence proxy \cite{schrijver2007}, and Lorentz-force
  components---consistent with established flare physics
  (Section~\ref{sec:results}).
\end{enumerate}

We are explicit about what is \emph{not} claimed: per-client SMOTE
contributes nothing on either substrate (zero synthetic samples on
the cleaned variant; an executed negative at natural prevalence); the
transport protocol moves full model weights, not data, and confers no
cryptographic privacy (the secure-aggregation sketch verifies
arithmetic cancellation only); the six clients are simulated, not real
agencies; and the downstream flare$\to$CME$\to$storm$\to$GIC chain
is discussed as a boundary, not evaluated.

The released framework also carries a documented revision history
(Appendix~\ref{sec:appendix}): a multi-version stability iteration in
which Dirichlet concentration, focal-loss alpha clamping, optimiser
choice, and dataset auditing were each adjusted in response to
diagnosed failure modes, and the audit findings reported here.
We keep that record because reproducible failure diagnosis is
transferable know-how for first FL deployments in the physical
sciences.

The remainder of the paper is organised as follows.
Section~\ref{sec:related} reviews flare prediction and federated
learning. Section~\ref{sec:problem} formalises the problem.
Section~\ref{sec:method} details the system architecture and
methodology. Section~\ref{sec:experiments} describes the experimental
setup, Section~\ref{sec:results} presents and analyses results, and
Section~\ref{sec:discussion} discusses operational integration and
engineering lessons. Section~\ref{sec:limitations} states limitations
and future work, and Section~\ref{sec:conclusion} concludes.
